\section*{الفصل \ref{ch:b3:atmospheric-dynamics} --- الديناميكا الجوية}

\begin{solution}{exo:b3:atmospheric-dynamics:1}
(أ)~$\zeta = 2\omega = 2\times2\pi/86\,400 = \qty{1.45e-4}{s^{-1}}$.
(ب)~$f = \num{7.29e-5}$ و\qty{1.26e-4}{s^{-1}}: تبلغ الدوّامة تقريبًا
ضعف \omterm{def:b3:atmospheric-dynamics:vorticity-kinds}{الدوّامية الكوكبية} عند \ang{30}~شمالًا، وأكثر قليلًا من
تلك عند \ang{60}~شمالًا.
\end{solution}

\begin{solution}{exo:b3:atmospheric-dynamics:2}
$\zeta = f_{30} - f_{40} = \num{7.29e-5} - \num{9.37e-5} =
\qty{-2.08e-5}{s^{-1}}$: مضادّة للإعصار.
\end{solution}

\begin{solution}{exo:b3:atmospheric-dynamics:3}
$k = 2\pi/(5\times10^6) = \qty{1.257e-6}{m^{-1}}$; $c = -\beta/k^2 =
-\num{1.62e-11}/\num{1.58e-12} = \qty{-10.3}{m/s}$: متجهة غربًا.
\end{solution}

\begin{solution}{exo:b3:atmospheric-dynamics:4}
$\beta = 2\Omega\cos\ang{60}/a = \qty{1.14e-11}{m^{-1}.s^{-1}}$;
$\lambda_s = 2\pi\sqrt{10/\num{1.14e-11}} = \qty{5870}{km}$؛ الدائرة،
\qty{20015}{km}، يتّسع لنحو $n_s = 3.4$ موجة.
\end{solution}

\begin{solution}{exo:b3:atmospheric-dynamics:5}
$q = f/h$ محفوظة: عبر السلسلة $\zeta = f(8/10) - f =
\qty{-2.06e-5}{s^{-1}}$، مضادّة للإعصار؛ وفوق الحوض $f(12/10) - f =
\qty{+2.06e-5}{s^{-1}}$، إعصارية. حيث تُمَطّ الأعمدة، في الجانب المحمي
من السلسلة وفوق الحوض تتولّد دوّامية إعصارية، وفيها
تميل المنخفضات الجوية إلى التكوّن.
\end{solution}

\begin{solution}{exo:b3:atmospheric-dynamics:6}
$k = 5/(a\cos\ang{45}) = \qty{1.11e-6}{m^{-1}}$, $\beta/k^2 =
\qty{13.1}{m/s}$: $c = \qty{1.9}{m/s}$، $c_g = \qty{28.1}{m/s}$. تقطع الطاقة
\qty{5000}{km} خلال \num{2.1}~يومًا، والنمط خلال \num{31}~يومًا.
\end{solution}

\begin{solution}{exo:b3:atmospheric-dynamics:7}
(أ)~$k^2 = \beta/U - l^2 = \num{1.079e-12} - \num{2.74e-13} =
\qty{8.05e-13}{m^{-2}}$, $k = \qty{8.97e-7}{m^{-1}}$: $\lambda =
\qty{7000}{km}$ و$n = 4.0$، بعدد موجات أقلّ مما عند $l = 0$. (ب)~$l^2 =
\num{1.097e-12} > \beta/U = \num{1.079e-12}$: $k^2$ سيكون سالبًا.
\end{solution}

\begin{solution}{exo:b3:atmospheric-dynamics:8}
$L_d = 0.012\times10^4/\num{1.031e-4} = \qty{1164}{km}$; $\sigma =
0.31\times\num{1.031e-4}\times\num{2e-3}/0.012 = \qty{5.3e-6}{s^{-1}}$، أي
زمن التضاؤل الأسّي $e$ البالغ \num{2.2}~يومًا؛ وطول الموجة $2\pi\times1164/1.61 =
\qty{4540}{km}$.
\end{solution}

\begin{solution}{exo:b3:atmospheric-dynamics:9}
$[v^*T^*] = (VT_0/2)\cos\delta$: \qty{25}{K.m/s}، \qty{12.5}{K.m/s} و0؛
دفق الحرارة \num{25100}، \num{12600} و\qty{0}{W/m^2}. عندما تكون
درجة الحرارة خارج الطور مع الريح بمقدار ربع طول موجة، فيتّجه الهواء الدافئ
جنوبًا بقدر ما يتّجه شمالًا: فلا يوجد نقل.
\end{solution}

\begin{solution}{exo:b3:atmospheric-dynamics:10}
بالتقريب الخطّي حول $U$ مع $\zeta' = \nabla^2\psi'$ و$v' =
\partial_x\psi'$: $(\partial_t + U\partial_x)\nabla^2\psi' + \beta\partial_x
\psi' = 0$. حيث $\psi' \propto \mathrm e^{\mathrm i(kx + ly - \omega t)}$،
$(\omega - Uk)K^2 = -\beta k$، ومنه $\omega = Uk - \beta k/K^2$. بالنسبة إلى
الهواء، تكون سرعة الطور $c - U = -\beta/K^2 < 0$ لكلّ موجة.
\end{solution}

\begin{solution}{exo:b3:atmospheric-dynamics:11}
$\varphi_H = [5\times9.81\times12\,000\times0.25/(3\times(\num{7.292e-5})^2
\times(\num{6.371e6})^2)]^{1/2} = \qty{0.477}{rad}$, \ang{27.3}.
$u_M = \Omega a\varphi_H^2 = \qty{106}{m/s}$؛ وبصورة دقيقة، $\Omega a\sin^2
\varphi_H/\cos\varphi_H = \qty{110}{m/s}$: ضعفين ونصفًا
من التيّار النفّاث المرصود.
\end{solution}

\begin{solution}{exo:b3:atmospheric-dynamics:12}
(أ)~عبر النطاق يُزاح السطح البيني بمقدار $\eta = sy$، $-L/2 \le y
\le L/2$، عن وضعه المستوي. طبقة ذات تباين كثافة مُخفَّض
تُزاح بمقدار $\eta$ تخزّن $\tfrac12\rho g'\eta^2$ لكلّ وحدة مساحة؛ ويُحسب المتوسّط
على النطاق، $\langle\eta^2\rangle = s^2L^2/12$. (ب)~$g' = 0.327$:
\num{1.5e5} و\qty{6.1e5}{J/m^2}، مقابل $\tfrac12\times0.7\times15^2
\times10^4 = \qty{7.9e5}{J/m^2}$ من طاقة الحركة: متقاربتان.
\end{solution}

\begin{solution}{pb:b3:atmospheric-dynamics:1}
\textbf{1.} $f_{45} = \qty{1.031e-4}{s^{-1}}$, $f_{50} =
\qty{1.117e-4}{s^{-1}}$; $\zeta = f_{45} - f_{50} = \qty{-8.6e-6}{s^{-1}}$.

\textbf{2.} $R = V/|\zeta| = 15/\num{8.6e-6} = \qty{1750}{km}$؛ في اتّجاه عقارب الساعة،
عائدةً نحو الجنوب.

\textbf{3.} $\zeta = f_{45}(9.5/10) - f_{50} = \qty{-1.38e-5}{s^{-1}}$، أكثر
مضادّةً للإعصار.

\textbf{4.} $\zeta = f_{45} - f_{40} = \qty{+9.4e-6}{s^{-1}}$، إعصارية: فهي
تعود شمالًا. وعند إزاحة الكتلة الهوائية في أيّ من الاتّجاهين، تكتسب الدوّامية
التي تعيدها: تأثيرٌ استعاديّ، وتذبذبٌ.

\textbf{5.} $\beta = \qty{1.62e-11}{m^{-1}.s^{-1}}$؛ تبلغ الدائرة
\qty{28306}{km} طولًا.

\textbf{6.} $n = 3$ إلى 8: أطوال الموجات 9435، 7076، 5661، 4718، 4044
و\qty{3538}{km}؛ $c = -21.5$، $-5.5$، 1.9، 5.9، 8.3 و\qty{9.9}{m/s}؛ $c_g
= 51.5$، 35.5، 28.1، 24.1، 21.7 و\qty{20.1}{m/s}.

\textbf{7.} $n = 3$ و4 تتحرّك غربًا، $n \ge 5$ شرقًا. تستغرق الموجة $n = 5$
$10^6/1.86 = \qty{5.4e5}{s}$، \num{6.2}~يومًا، لتتحرّك مسافة \qty{1000}{km}.

\textbf{8.} مع اتّجاه الريح، شرقًا: عند $c_g = \qty{24.1}{m/s}$ تقطع الطاقة
\qty{5000}{km} خلال \num{2.4}~يومًا.

\textbf{9.} $\lambda_s = 2\pi\sqrt{15/\num{1.62e-11}} = \qty{6050}{km}$,
$n_s = 4.7$.

\textbf{10.} $n_s = 5.7$ عند \qty{10}{m/s}، و4.1 عند \qty{20}{m/s}: ريح غربية
أقوى تفضّل موجات ثابتة أطول، وعددًا أقلّ حول
نصف الكرة.

\textbf{11.} $n_s = 5.8$ عند \ang{35}~شمالًا، و3.4 عند \ang{55}~شمالًا.

\textbf{12.} $\zeta = f(8/10) - f = \qty{-2.06e-5}{s^{-1}}$: الانضغاط فوق
الجبال يجعل العمود مرتفعًا جويًّا، والمطّ في الجانب الخلفي
منخفضًا جويًّا. تفرض الجبال نمطًا ثابتًا من الدوّامية على الجريان،
قوةً قسرية لا تتحرّك.

\textbf{13.} قوةٌ قسرية ثابتة تُنتج استجابةً مستقرّة؛ وتكون الاستجابة
أقوى ما يكون حين يكون النمط المقسور نفسه موجةً حرّةً تبقى
ساكنةً، أي الرنين عند $\lambda_s$.

\textbf{14.} $\Lambda = 25/10^4 = \qty{2.5e-3}{s^{-1}}$;
$|\partial T/\partial y| = \Lambda fT/g = \num{2.5e-3}\times\num{1.031e-4}
\times250/9.81 = \qty{6.6}{K}$ لكلّ \qty{1000}{km}.

\textbf{15.} $\mathrm{Ri} = 0.011^2/(\num{2.5e-3})^2 = 19$: أعلى بكثير من 0.25،
لا يوجد اضطراب ناتج عن القصّ.

\textbf{16.} $L_d = 0.011\times10^4/\num{1.031e-4} = \qty{1067}{km}$.

\textbf{17.} $2\pi\times1067/1.61 = \qty{4160}{km}$; $28306/4160 = 6.8$
موجات.

\textbf{18.} $2\pi\times1067/2.40 = \qty{2790}{km}$.

\textbf{19.} $\sigma = 0.31\times\num{1.031e-4}\times\num{2.5e-3}/0.011 =
\qty{7.3e-6}{s^{-1}}$، \num{0.63}~يوم$^{-1}$.

\textbf{20.} $\exp(3\times0.627) = 6.6$.

\textbf{21.} العاصفة، \qty{4160}{km}، أقصر من
\omterm{prop:b3:atmospheric-dynamics:stationary}{طول موجة روسبي الثابتة}، \qty{6050}{km}: بوصفها موجة روسبي تُحمل شرقًا بفعل
الرياح الغربية، وتتحرّك العواصف شرقًا عبر الموجات الطويلة
الأبطأ أو الواقفة.

\textbf{22.} تتجّه الريح غربًا مع الارتفاع: عندها فقط يكون الهواء المتحرّك نحو القطب
أدفأ من الهواء المتحرّك نحو خط الاستواء، فتنتقل الحرارة نحو القطب، وتنطلق
\omterm{def:b3:atmospheric-dynamics:ape}{طاقة الوضع المتاحة}.

\textbf{23.} الاضطراب الأسرع نموًّا يتزايد بعامل $e$ كلّ
\textbf{\num[text-series-to-math=true]{1.6}~يومًا} (\num{1.59}~يومًا حسب
الحساب).
\end{solution}
